\section{USAMO 2021}
        \begin{enumerate}
        	\item (\textit{USAMO 2021}) $\textbf{Note:}$ For any geometry problem whose statement begins with an asterisk $(*)$, the first page of the solution must be a large, in-scale, clearly labeled diagram. Failure to meet this requirement will result in an automatic 1-point deduction.


 $(*)$ Rectangles $BCC_1B_2,$ $CAA_1C_2,$ and $ABB_1A_2$ are erected outside an acute triangle $ABC.$ Suppose that
 \[\angle BC_1C+\angle CA_1A+\angle AB_1B=180^{\circ}.\]Prove that lines $B_1C_2,$ $C_1A_2,$ and $A_1B_2$ are concurrent.


 

	\item (\textit{USAMO 2021}) $\textbf{Note:}$ For any geometry problem whose statement begins with an asterisk $(*)$, the first page of the solution must be a large, in-scale, clearly labeled diagram. Failure to meet this requirement will result in an automatic 1-point deduction.


 The Planar National Park is a subset of the Euclidean plane consisting of several trails which meet at junctions. Every trail has its two endpoints at two different junctions whereas each junction is the endpoint of exactly three trails. Trails only intersect at junctions (in particular, trails only meet at endpoints). Finally, no trails begin and end at the same two junctions. 


 A visitor walks through the park as follows: she begins at a junction and starts walking along a trail. At the end of that first trail, she enters a junction and turns left. On the next junction she turns right, and so on, alternating left and right turns at each junction. She does this until she gets back to the junction where she started. What is the largest possible number of times she could have entered any junction during her walk, over all possible layouts of the park?


 

	\item (\textit{USAMO 2021}) $\textbf{Note:}$ For any geometry problem whose statement begins with an asterisk $(*)$, the first page of the solution must be a large, in-scale, clearly labeled diagram. Failure to meet this requirement will result in an automatic 1-point deduction.


 Let $n \geq 2$ be an integer. An $n \times n$ board is initially empty. Each minute, you may perform one of three moves:
If there is an L-shaped tromino region of three cells without stones on the board (see figure; rotations not allowed), you may place a stone in each of those cells.
If all cells in a column have a stone, you may remove all stones from that column.
If all cells in a row have a stone, you may remove all stones from that row.

                    \begin{figure}[H]
                        \centering
                        \includegraphics[width=0.5\textwidth]{Problem/USAMO_I_2021_3_Problem_Figure_1.png}
                    \end{figure}
                    
For which $n$ is it possible that, after some non-zero number of moves, the board has no stones?


 

	\item (\textit{USAMO 2021}) $\textbf{Note:}$ For any geometry problem whose statement begins with an asterisk $(*)$, the first page of the solution must be a large, in-scale, clearly labeled diagram. Failure to meet this requirement will result in an automatic 1-point deduction.


 A finite set $S$ of positive integers has the property that, for each $s \in S,$ and each positive integer divisor $d$ of $s$, there exists a unique element $t \in S$ satisfying $\text{gcd}(s, t) = d$. (The elements $s$ and $t$ could be equal.)


 Given this information, find all possible values for the number of elements of $S$.


 

	\item (\textit{USAMO 2021}) $\textbf{Note:}$ For any geometry problem whose statement begins with an asterisk $(*)$, the first page of the solution must be a large, in-scale, clearly labeled diagram. Failure to meet this requirement will result in an automatic 1-point deduction.


 Let $n \geq 4$ be an integer. Find all positive real solutions to the following system of $2n$ equations:


 
 \begin{align*} a_{1} &=\frac{1}{a_{2 n}}+\frac{1}{a_{2}}, & a_{2}&=a_{1}+a_{3}, \\ a_{3}&=\frac{1}{a_{2}}+\frac{1}{a_{4}}, & a_{4}&=a_{3}+a_{5}, \\ a_{5}&=\frac{1}{a_{4}}+\frac{1}{a_{6}}, & a_{6}&=a_{5}+a_{7} \\ &\vdots & &\vdots \\ a_{2 n-1}&=\frac{1}{a_{2 n-2}}+\frac{1}{a_{2 n}}, & a_{2 n}&=a_{2 n-1}+a_{1} \end{align*}



 

	\item (\textit{USAMO 2021}) $\textbf{Note:}$ For any geometry problem whose statement begins with an asterisk $(*)$, the first page of the solution must be a large, in-scale, clearly labeled diagram. Failure to meet this requirement will result in an automatic 1-point deduction.


 $(*)$ Let $ABCDEF$ be a convex hexagon satisfying $\overline{AB} \parallel \overline{DE}$, $\overline{BC} \parallel \overline{EF}$, $\overline{CD} \parallel \overline{FA}$, and

 \[AB \cdot DE = BC \cdot EF = CD \cdot FA.\]
Let $X$, $Y$, and $Z$ be the midpoints of $\overline{AD}$, $\overline{BE}$, and $\overline{CF}$. Prove that the circumcenter of $\triangle ACE$, the circumcenter of $\triangle BDF$, and the orthocenter of $\triangle XYZ$ are collinear.


 

\end{enumerate}